% 导言区
% \documentclass{article}
\documentclass{ctexart}

\usepackage{ctex}


% 正文区
\begin{document}
    \section{引言} % 构建小节
    西方称勾股定理为毕达哥拉斯定理，将勾股定理的发现归功于公元前 6 世纪的毕达哥拉斯学派。
    该学派得到了一个法则，可以求出可排成直角三角形三边的三元数组。毕达哥拉斯学派没有书面著作，该定理的严格表述和证明则见于欧几里德欧几里德，约公元前 330--275 年。
    《几何原本》的命题 47：“直角三角形斜边上的正方形等于两直角边上的两个正方形之和。”证明是用面积做的。
    我国《周髀算经》载商高（约公元前 12 世纪）答周公问

    西方称勾股定理为毕达哥拉斯定理，将勾股定理的发现归功于公元前 6 世纪的毕达哥拉斯学派。
    该学派得到了一个法则，可以求出可排成直角三角形三边的三元数组。\\毕达哥拉斯学派没有书面著作，该定理的严格表述和证明则见于欧几里德欧几里德，约公元前 330--275 年。
    《几何原本》的命题 47：“直角三角形斜边上的正方形等于两直角边上的两个正方形之和。”证明是用面积做的。
    我国《周髀算经》载商高（约公元前 12 世纪）答周公问

    西方称勾股定理为毕达哥拉斯定理，将勾股定理的发现归功于公元前 6 世纪的毕达哥拉斯学派。
    该学派得到了一个法则，可以求出可排成直角三角形三边的三元数组。\par 毕达哥拉斯学派没有书面著作，该定理的严格表述和证明则见于欧几里德欧几里德，约公元前 330--275 年。
    《几何原本》的命题 47：“直角三角形斜边上的正方形等于两直角边上的两个正方形之和。”证明是用面积做的。
    我国《周髀算经》载商高（约公元前 12 世纪）答周公问
    \section{实验方法}
    \section{实验结果}
    \subsection{数据} % 构建子小节
    \subsection{图表}
    \subsubsection{实验条件} % 构建下一级的小节
    \subsubsection{实验过程}
    \subsection{结果分析}
    \section{结论}
    \section{致谢}
\end{document}
